\section*{Hoofdstuk \ref{ch:g10:habitable-planet} --- Een bewoonbare planeet}

\begin{solution}{exo:g10:habitable-planet:1}
De temperatuur waarbij een perfecte straler precies het vermogen naar de ruimte
zou uitstralen dat de planeet van haar ster absorbeert. Het aardoppervlak is
warmer (\qty{288}{K} tegen ongeveer \qty{255}{K}) omdat broeikasgassen een deel
van het infrarood dat het uitzendt absorberen en een deel ervan terugsturen.
\end{solution}

\begin{solution}{exo:g10:habitable-planet:2}
Venus: gemeten \qty{737}{K} tegen een \omterm{def:g10:habitable-planet:equilibrium-temperature}{evenwichtstemperatuur} van ongeveer
\qty{227}{K}, een opwarming van ongeveer \qty{510}{K}.
\end{solution}

\begin{solution}{exo:g10:habitable-planet:3}
$S = 1361/1.52^2 \approx \qty{589}{W/m^2}$ (het gegevensblad geeft
\qty{586}{W/m^2}, met de exacte gemiddelde afstand).
\end{solution}

\begin{solution}{exo:g10:habitable-planet:4}
\qty{1}{bar}, \qty{330}{K}: vloeibaar (onder het kookpunt, ongeveer
\qty{373}{K}). \qty{0.001}{bar}, \qty{250}{K}: ijs (de druk ligt boven de
kromme tussen ijs en damp, die bij \qty{250}{K} op ongeveer \qty{0.0008}{bar} ligt).
\qty{100}{bar}, \qty{700}{K}: damp --- de temperatuur ligt boven de kritieke
temperatuur, dus er kan geen vloeistof bestaan.
\end{solution}

\begin{solution}{exo:g10:habitable-planet:5}
Een planeet die om een andere ster dan de Zon draait. Transits: de planeet
trekt vóór haar ster langs en maakt die bij elke omloop iets zwakker. Wiebelen:
de aantrekking van de planeet beweegt de ster heen en weer langs onze
gezichtslijn, waardoor de kleur van het sterlicht periodiek verschuift.
\end{solution}

\begin{solution}{exo:g10:habitable-planet:6}
$T = (1361 \times 0.706/(4 \times 5.67\times10^{-8}))^{1/4} \approx
\qty{255}{K} \approx \qty{-18}{\degreeCelsius}$. Met $A = 0$:
$\approx \qty{278}{K} \approx \qty{5}{\degreeCelsius}$.
\end{solution}

\begin{solution}{exo:g10:habitable-planet:7}
$\sqrt 4 = 2$: van $2 \times 0.99 = \qty{1.98}{au}$ tot $2 \times 1.70 =
\qty{3.40}{au}$. Op \qty{1}{au} zou de planeet vier keer de bestralingssterkte
van de Aarde ontvangen: ver binnen de binnenrand, te heet.
\end{solution}

\begin{solution}{exo:g10:habitable-planet:8}
$S = 1361/0.723^2 \approx \qty{2604}{W/m^2}$;
$T_{\mathrm{eq}} = (2604 \times 0.23/(4\sigma))^{1/4} \approx \qty{227}{K}$.
Het oppervlak is ongeveer \qty{510}{K} heter: de dikke atmosfeer van
koolstofdioxide (\qty{92}{bar}) veroorzaakt een enorm broeikaseffect.
\end{solution}

\begin{solution}{exo:g10:habitable-planet:9}
Bij \qty{6.36}{mbar} is water alleen vloeibaar tussen ongeveer
\qty{0}{\degreeCelsius} en \qty{0.5}{\degreeCelsius}. Bij \qty{10}{\degreeCelsius}
ligt de druk onder de kookkromme: het water kookt weg tot damp; 's nachts bevriest het.
\end{solution}

\begin{solution}{exo:g10:habitable-planet:10}
$11.2/2.4 \approx 4.7$. Een gasmolecuul moet bijna vijf keer sneller bewegen
om aan de Aarde te ontsnappen dan aan de Maan; in miljarden jaren liet de
zwakke zwaartekracht van de Maan haar gassen ontsnappen, terwijl die van de
Aarde het grootste deel ervan vasthield.
\end{solution}

\begin{solution}{exo:g10:habitable-planet:11}
$T = (1361 \times 0.50/(4\sigma))^{1/4} \approx \qty{234}{K}$, ongeveer
\qty{-39}{\degreeCelsius}. Het ijs zou zoveel zonlicht weerkaatsen dat de
planeet koud genoeg zou blijven om het te behouden: een zichzelf versterkende kringloop.
\end{solution}

\begin{solution}{exo:g10:habitable-planet:12}
$T_{\mathrm{eq}} = (S_0(1-A)/(4\sigma d^2))^{1/4} = (S_0/(4\sigma))^{1/4}
(1-A)^{1/4} (d^2)^{-1/4}$; $(1361/(4\sigma))^{1/4} \approx \qty{278.3}{K}$ en
$(d^2)^{1/4} = \sqrt d$. Vier keer zo ver: $T$ gedeeld door $\sqrt 4 = 2$.
\end{solution}

\begin{solution}{exo:g10:habitable-planet:13}
$\sqrt d = 278.3 \times 0.706^{1/4}/273 \approx 0.934$, dus $d \approx
\qty{0.873}{au}$, binnen de binnenrand (\qty{0.99}{au}). De
\omterm{def:g10:habitable-planet:equilibrium-temperature}{evenwichtstemperatuur} houdt geen rekening met het broeikaseffect: een echte
planeet met een atmosfeer is warmer dan $T_{\mathrm{eq}}$, zodat ze al verder
naar buiten, op ongeveer \qty{0.99}{au}, oververhit raakt en haar water
verliest. De randen komen uit modellen van de hele atmosfeer, niet uit
$T_{\mathrm{eq}}$ alleen.
\end{solution}

\begin{solution}{exo:g10:habitable-planet:14}
Tussen het smeltpunt, ongeveer \qty{0}{\degreeCelsius}, en het kookpunt,
ongeveer \qty{7}{\degreeCelsius}. Bij \qty{15}{\degreeCelsius} zou het meer
koken.
\end{solution}

\begin{solution}{exo:g10:habitable-planet:15}
$\sqrt{0.5} \approx 0.707$: van $0.707 \times 0.99 \approx \qty{0.70}{au}$
tot $0.707 \times 1.70 \approx \qty{1.20}{au}$. De planeet ligt precies op de
buitenrand. De randen zijn modeluitkomsten en de schaling is bij benadering,
dus het valt niet te zeggen of ze erbinnen of erbuiten ligt.
\end{solution}

\begin{solution}{pb:g10:habitable-planet:1}
\textbf{1.} $0.00151/0.0485^2 \approx 0.642$.

\textbf{2.} $0.642 \times 1361 \approx \qty{874}{W/m^2}$.

\textbf{3.} Ertussenin: ongeveer drie keer minder dan bij Venus, anderhalf keer
meer dan bij Mars.

\textbf{4.} $365.25/11.18 \approx 32.7$.

\textbf{5.} $(874/(4 \times 5.67\times10^{-8}))^{1/4} \approx \qty{249}{K}$.

\textbf{6.} $(874 \times 0.23/(4\sigma))^{1/4} \approx \qty{173}{K}$.

\textbf{7.} Geen van beide: \qty{249}{K} ligt nog onder \qty{273}{K}. Een
heldere laag wolken of ijs verlaagt de \omterm{def:g10:habitable-planet:equilibrium-temperature}{evenwichtstemperatuur} met meer dan
\qty{75}{K}: wat een planeet bedekt, telt even zwaar als haar afstand.

\textbf{8.} $\sqrt{0.00151} \approx 0.0389$.

\textbf{9.} $0.99 \times 0.0389 \approx \qty{0.0385}{au}$.

\textbf{10.} $1.70 \times 0.0389 \approx \qty{0.066}{au}$.

\textbf{11.} Ja: $0.0385 < 0.0485 < 0.066$.

\textbf{12.} $d_\odot = 0.0485/0.0389 \approx \qty{1.25}{au}$: in de binnenste
helft van de \omterm{def:g10:habitable-planet:habitable-zone}{bewoonbare zone} van de Zon, tussen de Aarde en Mars.

\textbf{13.} $1.055^{1/3} \approx 1.018$.

\textbf{14.} $11.2 \times \sqrt{1.055/1.018} \approx \qty{11.4}{km/s}$, ongeveer
die van de Aarde: haar zwaartekracht zou een atmosfeer kunnen vasthouden (of ze
er een heeft, hangt af van haar geschiedenis).

\textbf{15.} Het wiebelen van de ster verraadt alleen het deel van haar
beweging langs onze gezichtslijn; omdat de helling van de baan onbekend is, is
de gevonden massa alleen de werkelijke massa als we de baan precies van opzij
zien, en anders een minimum.

\textbf{16.} De druk van het tripelpunt, \qty{611.657}{Pa}, ongeveer
\qty{0.006}{bar}.

\textbf{17.} $(874 \times 0.706/(4\sigma))^{1/4} \approx \qty{228}{K}$.

\textbf{18.} $288 - 228 = \qty{60}{K}$: Proxima~b zou een opwarming door het
broeikaseffect van \textbf{ongeveer \qty{60}{K}} nodig hebben, bijna twee keer
de \qty{33}{K} van de Aarde.
\end{solution}
